\documentclass[11pt]{article}
\usepackage[margin=1in]{geometry}
\usepackage[T1]{fontenc}
\usepackage{lmodern,microtype,amsmath,amssymb,amsthm,booktabs,graphicx,enumitem}
\usepackage[colorlinks=true,linkcolor=blue,urlcolor=blue]{hyperref}
\setlength{\parindent}{0pt}
\setlength{\parskip}{4pt}
\setlist[itemize]{leftmargin=*,itemsep=3pt,topsep=3pt}
\newtheorem{proposition}{Proposition}
\newcommand{\dd}{\mathrm d}
\newcommand{\pageheading}[1]{\section*{#1}}
\begin{document}
\begin{center}
{\LARGE Housing allocation across generations}\par
\vspace{4pt}
A simple argument and its extensions\par
\smallskip
Discussion note for Tommaso De Santo \quad September 5, 2026
\end{center}

\pageheading{1. The question}

Can borrowing limits and restricted rental housing leave too much of the
existing housing stock with older households? The useful claim is that a
young household may value additional space more than an older owner values
keeping it. A small reallocation can then improve welfare while leaving
fertility unchanged. We can establish this before discussing a policy.

\textbf{Households and preferences.} Take an equilibrium of the existing
two-period OLG model. At each date there are young and old households, with
an infinite succession of cohorts. Young households choose consumption,
housing, completed fertility, saving, and tenure. Their flow utility and the
old household's flow utility remain:
\begin{align}
u_t^Y(c,h,n)&=\log(c-\chi n)+\alpha\log(h-\kappa n)+\vartheta_t\log n,\\
u^2(c^2,h^2,e)&=\log c^2+\gamma\log h^2+\omega_B\log e.
\end{align}
Each child uses $\chi$ goods and $\kappa$ housing; $e$ is the warm-glow estate.
Old-age utility is discounted by $\beta$; ownership adds the taste $\xi_i$.
The original lifetime value functions $W_t^m$ and old-age value functions
$V_t^m$ retain their definitions. Write $x=c-\chi n$ and $s=h-\kappa n$
for the goods and space available to adults.

\textbf{Housing and finance.} The fixed housing stock is $\bar H$. Rental
homes satisfy $h\leq h_R^{\max}$, and owner homes satisfy
$h\leq h_O^{\max}$, with $h_R^{\max}<h_O^{\max}$. A mortgage finances the
share $\phi$ of the purchase. A young buyer with liquid wealth $b_i$ must meet:
\begin{equation}
 (1-\phi)P_t h\leq b_i.
 \label{eq:simple_purchase}
\end{equation}
Only liquid wealth is available at closing. Income and ordinary rebates arrive
later. Old owners can sell part of their homes, but cannot buy more or let
their retained housing. The original saving, retention, and estate restrictions
remain part of the household problems.

Households save in a world bond priced at $q$. There is no capital-gains
tax in this core. The ordinary property-tax rate is $\tau^p$, so the current value of
one unit of rent is:
\begin{equation}
 u_t=q r_t=(1+q\tau^p)P_t-qP_{t+1}.
 \label{eq:simple_user_cost}
\end{equation}
This is the cost of using a unit of housing for the period, including its
subsequent resale value.

\textbf{The comparison.} Hold each household's fertility and tenure fixed.
The direct planner may relax private financing limits, as already agreed,
while respecting the housing stock, physical tenure restrictions, and real
resources. We compare each household's lifetime welfare, including the initial
old and all future cohorts. The construction below leaves every future real
allocation unchanged.

\newpage
\pageheading{2. The short allocation argument}

Consider a young owner $i$ and an old owner $j$. Their willingness to give up
current consumption for an additional housing unit is:
\begin{equation}
 MV_i^Y=\frac{\alpha x_i}{s_i},\qquad
 MV_j^O=\frac{\gamma c_j^2}{h_j^2}.
 \label{eq:simple_values}
\end{equation}
These are marginal rates of substitution, expressed in the same goods units.
They do not compare raw marginal utilities across people.

\begin{proposition}[Housing misallocation]
Suppose a positive mass of young owners can physically occupy more housing
and a positive mass of old owners can release some. If, on matched groups,
$MV_i^Y>MV_j^O$, a sufficiently small owner-to-owner reallocation can make
young households better off and leave everyone else equally well off, holding
fertility fixed. The benchmark admits direct goods compensation and financial
settlement while relaxing private financing restrictions. Reallocation has no
additional real cost.
\end{proposition}

\textbf{Proof.} Select equal-mass groups with a strict value gap and room for
a small change. Transfer $\epsilon$ housing per household from old to young.
The old household requires the following extra consumption to retain its utility:
\begin{equation}
 D_j(\epsilon)=c_j^2\left[
 \left(\frac{h_j^2}{h_j^2-\epsilon}\right)^\gamma-1\right].
 \label{eq:simple_compensation}
\end{equation}
Give it this amount from the young household's current consumption. Goods and
housing balance exactly. The young household's utility change is
$\log(1-D_j(\epsilon)/x_i)+\alpha\log(1+\epsilon/s_i)$; its derivative at
zero is $(MV_i^Y-MV_j^O)/x_i>0$. Small changes preserve positive consumption
and physical feasibility. Estates, future consumption and occupancy, and all
other households' allocations stay fixed. Title settlement replaces the old
seller's lost estate title with an equal bond claim; the young buyer's later
sale of the extra title offsets its financing. Existing creditors and the
aggregate external-asset position are unchanged.\hfill$\square$

\begin{center}
\includegraphics[width=.87\textwidth]{output/model/simplified_olg_amendments/simple_allocation_figure.pdf}
\end{center}
{\small\textbf{Figure 1.} A compensated reallocation in Appendix A's
equilibrium, holding old utility fixed. Shading marks the value gap; its area
is not a welfare measure. Equal values maximize the young household's gain
along this comparison, without solving the economy's full first best.}

\newpage
\pageheading{3. Why the constraints matter}

\textbf{The borrowing limit creates the value gap.} In the case we have been
studying, young owners save a positive amount, their physical housing limit is
slack, and their old-age retention and estate constraints are slack. If the
purchase limit strictly restricts housing, their optimal choices imply
$MV_i^Y>u_t$. An old owner with slack retention and estate constraints has
$MV_j^O=u_t$. The equilibrium therefore supplies the strict inequality used
in Proposition 1. Appendix A constructs mixed-tenure steady states satisfying
these conditions; they are not empty assumptions.

The proposition itself needs only the actual value gap and feasible
compensation. If another household constraint binds, compare the two values
in \eqref{eq:simple_values} directly. The special equality to user cost may
then fail.

\textbf{The rental restriction limits an alternative way to obtain space.}
A young family may be unable to rent the home it wants, and unable to fund the
down payment on a larger owned home. This is a useful account of housing access.
The local proof, however, transfers housing between owners and leaves rental
housing unchanged. A renter at a physical size cap cannot receive more space
under the same tenure when the planner must respect that cap. Its high housing
value alone is therefore insufficient to prove misallocation.

\textbf{An ownership taste matters for this interpretation.} The original
model gives households a taste $\xi_i$ for ownership. With a sufficiently
positive taste, a household may prefer constrained ownership even if larger
rentals are available. Indeed, Appendix A gives an equilibrium with slack
rental limits and a strict owner value gap. Removing the redundant rental
limit leaves that equilibrium unchanged. Both constraints need not bind for
the owner-to-owner inefficiency to exist.

There is a useful elementary comparison if we want both access restrictions
to be essential. Suppose an owner receives no ownership-taste benefit, and
the rental menu contains its current and old-age housing choices. Its entire
owner allocation can be reproduced by renting and choosing financial saving:
\begin{equation}
 a_R'=a_O'+qP_{t+1}h-\phi P_t h.
 \label{eq:simple_replication}
\end{equation}
Substituting \eqref{eq:simple_user_cost} into the original budgets verifies
both dated budgets exactly. If this saving is nonnegative and the current
rental cap leaves room, a household with $MV_i^Y>u_t$ can then improve by
renting a little more space. Such a household would not choose constrained
ownership without an ownership benefit. Rental segmentation can prevent that
alternative. Removing the ownership taste would make this joint role sharper,
but would be a separate model choice, not a notation change.

\textbf{What OLG adds.} Young households need purchase cash while old owners
already hold titles. Housing needs and financial positions differ over life,
and the same stock passes through successive cohorts. The local proof remains
a simple exchange argument with fixed continuation allocations. It establishes
housing misallocation without invoking capital overaccumulation or benefits
from additional births.

\newpage
\pageheading{4. A separate fertility prediction}

More housing reduces the space cost of children, but paying for it leaves
less for their goods costs. These effects can be separated in the original
household equations.

Take the conditional household problem with positive saving and unchanged
old-age constraints. At a binding housing limit, it gives:
\begin{equation}
 (1+k)x+\chi n+uh=\widetilde w,\qquad
 \frac{\vartheta}{n}=\frac{\chi}{x}+\frac{\alpha\kappa}{h-\kappa n}.
 \label{eq:simple_fertility_system}
\end{equation}
Here $k=\beta(1+\gamma+\omega_B)$ with interior old choices, and
$\widetilde w=y+b+T_t+qT_{t+1}$. If old renter housing is capped, instead use
$k=\beta(1+\omega_B)$ and subtract the discounted fixed old rent from resources.

Let $p$ be the net lifetime payment per marginal unit of current housing.
Relaxing only its cap gives $p=u_t$ at fixed prices and resources. If both
binding rental caps increase equally, $p=u_t+qu_{t+1}$. Accompanying transfers
also enter $p$. The budget gives $(1+k)\dd x=-\chi\dd n-p\dd h$.
Substituting this into the fertility condition yields:
\begin{equation}
 \frac{\dd n}{\dd h}=
 \frac{\displaystyle\frac{\alpha\kappa}{s^2}
 -\frac{\chi p}{(1+k)x^2}}
 {\displaystyle\frac{\vartheta}{n^2}
 +\frac{\chi^2}{(1+k)x^2}+\frac{\alpha\kappa^2}{s^2}}.
 \label{eq:simple_fertility_derivative}
\end{equation}
The denominator is positive. Fertility rises precisely when the space benefit
in the numerator exceeds the resource cost. There is no borrowing multiplier
in this condition.

\textbf{A sufficient restriction in primitives.} With stationary prices,
zero property tax, interior old choices, and a strictly binding down-payment
limit, the existing sufficient restriction is:
\begin{equation}
 \frac{(1-q)b}{(1-\phi)(y+b)}
 \geq\frac{\alpha}{1+\beta(1+\gamma+\omega_B)+\alpha}.
 \label{eq:simple_primitive}
\end{equation}
It implies that a marginal increase in housing raises fertility for every
$p<MV^Y$ while the stated constraints continue to bind. The left side is
housing expenditure at the purchase limit relative to lifetime resources;
the price cancels because the limit is inversely proportional to price.
This is sufficient, not necessary, and the zero-property-tax specialization
is explicit. Endogenous rebates cannot generally be treated as primitives.

\textbf{Why the qualification matters.} For the constructed household values
$1+k=2$, $\alpha=\chi=x=s=n=1$, $\kappa=0.1$, $\vartheta=1.1$, and
$u=p=0.5$, more housing improves welfare but
$\dd n/\dd h=-15/161$. At $p=0$ the derivative is $10/161$.
The difference is the payment for the extra space.

This is a conditional household prediction. The fixed-fertility allocation
proof holds future choices fixed; this calculation allows optimal saving and
old-age choices to adjust. They are distinct comparisons. A policy also changes
prices, tenure, and later cohort sizes, which must be included before claiming
an aggregate fertility increase.

\newpage
\pageheading{5. Population levels and the second figure}

Let $Y_t$ and $O_t$ count young and old households, and let $\bar n_t$ be
mean young fertility. The parameter $\nu$ converts births into entering
households. From $Y_{t+1}=\nu\bar n_tY_t$, two paths with the same initial
young cohort satisfy:
\begin{equation}
 \frac{Y_t^1}{Y_t^0}=\prod_{s=0}^{t-1}
 \frac{\bar n_s^1}{\bar n_s^0},\qquad O_t^j=Y_{t-1}^j.
 \label{eq:simple_population_product}
\end{equation}
Temporary fertility differences can therefore leave permanent population
differences. The old cohort follows one period later. What matters is the
cumulative fertility difference, which can include some negative terms.

At a positive steady state, $\bar n^*=1/\nu$ and $Y^*=O^*$. Housing clearing
therefore fixes total household population:
\begin{equation}
 N_{\rm hh}^*=
 \frac{2\bar H}{\bar h^{Y*}+\bar h^{O*}}.
 \label{eq:simple_population_stock}
\end{equation}
With a fixed stock, a larger population requires less housing per household
on average. Young housing may rise if housing used by older households falls
enough. Replacement fertility alone does not determine the population level.

\textbf{A simple supply extension.} Multiplying the exogenous housing stock
by $a>0$ gives a corresponding steady state with every cohort mass multiplied
by $a$ and unchanged prices, individual choices, and rebates per household.
Housing supply and demand scale together. The rebate is unchanged because
property-tax revenue and recipients also scale together. Household problems
remain identical with the entrant distribution and world bond price fixed.
Thus stationary population scales with the stock. This uses the existing
closure: the extra stock adds no further entrant-income or fiscal channel.

\begin{center}
\includegraphics[width=\textwidth]{output/model/simplified_olg_amendments/simple_population_figure.pdf}
\end{center}
{\small\textbf{Figure 2.} A prescribed fertility sequence and the population
path implied by \eqref{eq:simple_population_product}. Fertility returns to
replacement with more households. These curves do not establish an equilibrium
transition or its speed.}

We still need to establish an equilibrium path between the relevant steady
states. The existing continuation argument assumes regularity along the whole
reform; simple primitive conditions remain unavailable. The stationary ranking
also applies to persons when adults per household and the fixed timing and
residence rules for counting children are unchanged. During a transition,
resident children must also be counted.

\newpage
\pageheading{6. Constrained efficiency and public finance}

\textbf{The basic result already identifies misallocation.} Proposition 1
exhibits a feasible allocation that improves welfare. It respects physical
housing restrictions and relaxes private finance. It neither characterizes
the full first best nor requires a government lending policy.

\textbf{A constrained comparison asks an additional question.} Can an
authority improve welfare using a specified set of instruments while
households continue to choose at clearing prices? The restrictions on the
authority must be stated before answering. Retaining the private mortgage
formula while allowing public credit is a different comparison from retaining
the household's effective financing opportunities. Bisin's lecture notes,
section 3.2.4, make the choice of constrained feasible allocations explicit.

The current research contains two results with different scopes:
\begin{itemize}
\item A complete local improving path under committed household-specific
  payments when young and old, enforceable future taxes, and explicit initial
  title ownership with access to finance. Fertility is fixed individually.
  Every market clears, but the path returns to the original steady state.
\item An obstruction to one initial round of transfers near one studied
  stationary regime, for nearby converging paths on the maintained household
  branches. This does not establish efficiency under every other transfer class.
\end{itemize}

These results remain useful background. The second does not show that public
loans are the only possible remedy. The first supplies financing across a
household's life and should be described accordingly. Neither is required
for the short allocation argument.

\textbf{The loan interpretation is a later extension.} For later cohorts,
the committed-transfer construction advances purchase cash and collects
repayment plus part of the household's housing gain. A government with access
to world bonds could potentially replace the separate financier. All original
property claims and promised estate payments would still need to be honored.
A voluntary loan programme also changes participation and tenure choices;
its proof is not obtained just by relabeling identity-based transfers.

There are established precedents for the broad financing role. Auerbach and
Kotlikoff's compensation authority balances in present value; their payments
to future households occur on entry. Schoonbroodt and Tertilt use public
transfers to overcome parent-child transfer restrictions. With binding
liquidity constraints, the timing and eligibility of our proposed payments
remain substantive adaptations.

\textbf{Recommendation.} Use ``housing misallocation'' for the core result.
Retain constrained efficiency as an optional strengthening under a precisely
defined benchmark. Discuss public financing after deciding whether that
strengthening adds enough to justify its assumptions and length.

{\small
\textbf{References for these distinctions:}
\href{https://bpb-us-e1.wpmucdn.com/wp.nyu.edu/dist/c/16384/files/2019/11/Lecture-notes-Oct2014.pdf}{Bisin (2014), lecture notes, section 3.2.4};
\href{https://kotlikoff.net/wp-content/uploads/2019/04/Dynamic-Fiscal-Policy.pdf}{Auerbach and Kotlikoff (1987), chapter 5, B.3, pp. 62--64};
\href{https://tertilt.vwl.uni-mannheim.de/research/JET2013.pdf}{Schoonbroodt and Tertilt (2014), sections 5.2--5.3}.
These sources motivate the comparisons; our housing claims require the
model-specific proofs.}

\newpage
\pageheading{7. Assessment and opportunities}

\textbf{The core is ready to be stated simply.} A financial restriction can
prevent a young household from obtaining housing it values more than an old
owner does. A direct reallocation with consumption compensation establishes
the welfare content. The proof takes a paragraph. Its role is to define the
misallocation and explain the first figure. There is no need to present it as
an unusual welfare theorem.

\textbf{The useful clarification concerns tenure.} Rental segmentation helps
explain why a household cannot obtain its preferred space by renting. Yet the
ownership taste allows the borrowing distortion to persist without a binding
rental cap. The existing model supports the statement that the two restrictions
can coexist with misallocation. It does not support the claim that both are
necessary. The replication argument on page 3 shows how removing the ownership
taste could sharpen their joint role, at the cost of a model change.

\textbf{My recommendation is to retain the present model for the first pass.}
Preserve its notation and household problems, state the conditional allocation
result, and explain the two constraints accurately. Do not change preferences
simply to make both restrictions necessary. If their joint necessity matters
to the paper's message, consider the version without ownership tastes as an
explicit alternative.

The extensions have different value and cost:
\begin{itemize}
\item \textbf{Real reallocation costs: keep compact.} If
  $C(0)=0$ and $C'(0)=K$, replace the condition by
  $MV_i^Y>MV_j^O+K$. A positive fixed moving cost requires a finite transfer;
  the exact compensation formula supplies the test. Compensation paid to the
  seller is not itself a real resource cost.
\item \textbf{Fertility: retain as a separate conditional prediction.}
  Equation \eqref{eq:simple_fertility_derivative} expresses the space benefit
  and resource cost directly. The primitive restriction is available as a
  sufficient specialization, with its limitations stated once.
\item \textbf{Population: retain the identities and the stationary scaling
  result.} They explain how replacement fertility can coexist with different
  population levels. A general analytical transition after a credit reform
  remains a further result; a schematic figure cannot establish it.
\item \textbf{Constrained efficiency: optional strengthening.} Its value
  depends on a natural and clearly stated instrument class. The current
  committed-finance proof can remain outside the short core while that choice
  is considered.
\item \textbf{Public loans, endogenous rental supply, and gains taxation:
  later branches.} Each asks an additional economic question. Their detailed
  policy design would make the initial illustration harder to read.
\end{itemize}

For the paper, the core could occupy two or three pages after the household
setup, including Figure 1. The fertility prediction and demographic comparison
can follow briefly; longer finance and transition arguments belong in an
appendix only if the paper uses them.

\textbf{Proposed order for discussion.} First settle the allocation benchmark
and the wording of Proposition 1. Then choose how much of the fertility and
population extensions to include. Finally consider whether the constrained
comparison adds enough to keep. Revisit ownership tastes only if the paper's
message requires both access restrictions to be essential. No change to those
preferences is adopted here.

\newpage
\pageheading{Appendix A. The conditions can hold in equilibrium}

The following construction checks that Proposition 1 can apply to positive
masses of households in a mixed-tenure steady state. It chooses primitives
to support an allocation directly. Income and the ownership-taste location
vary across the family; this is not a tax comparative static or a calibration.

Set $P=1$, $q=0.5$, $\phi=0.8$, $b=0.2$, $\bar H=2$, and $\nu=2$.
Use $\alpha=\beta=\omega_B=0.4$, $\gamma=0.3$, $\chi=0.15$, $\kappa=0.5$,
$h_R^{\max}=0.25$, and $h_O^{\max}=2$. Choose young-owner adult consumption,
housing, and fertility as $(x,h,n)=(1,1,0.75)$. The fertility condition sets
$\vartheta=0.3525$. For any $0\leq\tau^p<0.28$, user cost is
$u=(1+\tau^p)/2<0.64=MV^Y$. Old-owner choices are:
\begin{equation}
 c^2=0.8,\qquad h^2=0.24/u,\qquad e=0.64.
 \label{eq:simple_witness_old}
\end{equation}
Retention and estate bounds are strictly slack throughout this interval.

Let $w=1.7925+u$. With both rental caps binding, the renter's fertility is the
unique root in $(0,0.5)$ of:
\begin{equation}
 \frac{0.3525}{n^R}=\frac{0.15}{x^R}
 +\frac{0.2}{0.25-0.5n^R},\qquad
 x^R=\frac{1.7925+0.625u-0.15n^R}{1.56}.
 \label{eq:simple_witness_renter}
\end{equation}
The left side decreases from infinity and the right side increases to
infinity. Their root has $x^R>1.30$, which makes both rental caps strictly
binding under the original housing conditions. Both tenures save positively.

Set the owner share to $\pi^O=(0.5-n^R)/(0.75-n^R)\in(0,1)$, so average
fertility is replacement. Define housing demand per young-old pair as
$d_h=\pi^O(1+0.24/u)+(1-\pi^O)0.5$. Choose
$Y=O=2/d_h$, $T=q\tau^p d_h/2$, and $y=w-b-(1+q)T>0$.
Housing and the property-tax budget then clear exactly. With any taste scale
$\sigma_\xi>0$, choose
$\bar\xi=\sigma_\xi\log[\pi^O/(1-\pi^O)]-W^O+W^R$ to support the owner
share under the original logistic rule. The dated budgets and first-order
conditions give household optima because their feasible sets are convex and
their utility is concave.

\textbf{A rental cap need not bind.} A second construction sets
$\tau^p=0$, $h_R^{\max}=1.5$, and owner fertility $n^O=0.49$, retaining
$(x,h)=(1,1)$ and the other displayed primitives. Set
$\vartheta=61397/302000$ and $y=2.0535$. Unconstrained renter demands and the
same tenure-share construction give:
\begin{center}
\begin{tabular}{@{}lrr@{}}
\toprule
 & Both rental caps bind & Both rental caps slack\\
\midrule
Property tax & $0.05$ & $0$\\
Owner share & $0.523810$ & $0.139390$\\
Young renter housing & $0.250000$ & $1.040368$\\
Young owner housing value & $0.640000$ & $0.529801$\\
Old owner housing value & $0.525000$ & $0.500000$\\
\bottomrule
\end{tabular}
\end{center}
The second equilibrium remains an equilibrium if its redundant rental limit
is removed. These are two different constructed economies, not a comparison
holding all primitives fixed while changing the rental cap.

\end{document}
